% Wearable-defined exercise counterfactual and advisory arc.
% Values are sourced from outputs/study2_exercise_stage0*,
% outputs/study2_exercise_path_b*, and outputs/study2_exercise_advisory/.

\section{From Residual Audit to Exercise Counterfactual: Identifiability, an Imposed Planning Response, and a Deployable Advisory}
\label{sec:exercise_counterfactual}

\subsection{Why a Counterfactual, and the Identifiability Question}
\label{subsec:ex_cf_question}

The frozen-model residual audit found no exercise-specific MAE penalty, but it
did find a horizon-growing positive signed residual in moderate and vigorous
wearable-defined episodes. Under the notebook convention
$b_h=\hat q_{50,t}^{(h)}-G_{t+h}$, this means that the model over-predicts the
later glucose trajectory and under-represents the exercise-associated drop.
This provisional residual headroom motivates a natural question: what would
glucose do if this person exercised?

That question is not answered by detector validation. Detectability and
identifiability are distinct. Detection is a signal-processing problem:
steps and HR can locate an ambulatory exercise-like state. Identification is
a confounding-and-overlap problem: comparable non-exercise outcomes must exist
for the same glucose state, time, and participant context. A detector can be
accurate while the magnitude of the associated counterfactual response remains
non-identifiable.

\subsection{Stage 0: The Magnitude Is Not Identifiable on AI-READI}
\label{subsec:ex_stage0}

Stage~0 attempted matched-control de-confounding. For each detected bout, the
mean-reversion drift in matched non-exercise windows was subtracted from the
observed 60-minute change. Cross-patient matching produced numerical overlap
on baseline glucose $g_0$ and time of day, but not participant comparability.
With the $\pm10$~mg/dL glucose caliper, $92.23\%$ to $97.24\%$ of controls were
drawn from another participant across the five $g_0$ bins. In the
$70$--$100$~mg/dL bin, the apparent effect was significantly positive,
$+2.50$~mg/dL with 95\% CI $[+0.73,+4.31]$. Because the contrast is largely
between people, this is a signature of residual between-person confounding,
not evidence of a physiological rise.

Within-patient matching removed that low-$g_0$ result. The corresponding
estimates were $+0.33$~mg/dL $[-5.40,+6.19]$ in the $70$--$100$ bin and
$-0.10$~mg/dL $[-3.53,+3.23]$ in the $100$--$130$ bin. The cross-patient
low-$g_0$ result therefore does not survive within-patient matching. This is
a selection-aware statement, not evidence of a zero effect. Comparability
came at the cost of positivity: only $48$ of $510$ low-$g_0$ bouts were
retained ($9.4\%$), and the two highest-$g_0$ bins retained only $17$ and $7$
bouts. Retention in the $200$--$300$~mg/dL bin was $7.1\%$.

\begin{table}[H]
  \centering
  \caption{%
    \textbf{Stage 0 matching audit.}
    Cross-patient matching supplies overlap but relies almost entirely on
    other participants. Within-patient matching supplies comparability but
    loses positivity. The low-$g_0$ positive contrast does not survive the
    within-patient restriction.%
  }
  \label{tab:exercise_stage0}
  \small
  \setlength{\tabcolsep}{3pt}
  \begin{tabular}{lllll}
    \toprule
    Scheme and $g_0$ bin & Caliper & Controls & Cross-patient & Effect at 60 min, mg/dL \\
    \midrule
    Cross, 70--100   & $\pm5$  & 15{,}141 & 97.60\% & $+3.78$ $[+2.09,+5.62]$ \\
    Cross, 70--100   & $\pm10$ & 15{,}300 & 96.12\% & $+2.50$ $[+0.73,+4.31]$ \\
    Within, 70--100  & $\pm5$  & 364       & 0\%     & $+0.33$ $[-5.40,+6.19]$ \\
    Within, 100--130 & $\pm5$  & 1{,}272   & 0\%     & $-0.10$ $[-3.53,+3.23]$ \\
    Within, 160--200 & $\pm5$  & 82        & 0\%     & $+9.91$ $[-16.95,+37.49]$ \\
    Within, 200--300 & $\pm5$  & 44        & 0\%     & $-31.45$ $[-63.24,-2.26]$ \\
    \bottomrule
  \end{tabular}
\end{table}

The within-patient curve is non-monotonic and changes sign across baseline
bins. Thus no matching scheme supplies participant comparability and adequate
positivity simultaneously for this estimand. The strong Stage~0 finding is
that the exercise-response \emph{magnitude is not identifiable on AI-READI}.
The de-confounded curve is not a defensible calibration target.

\subsection{Path B: A Structured Exercise Head with Imposed Magnitude}
\label{subsec:ex_pathb}

Path~B retains what AI-READI can support without treating the Stage~0 curve as
identified. The magnitude is imposed from the descriptive, associational
moderate-strain kernel of $16.66$~mg/dL, while AI-READI supplies timing, shape,
and live glucose-state gating. The output is an imposed planning response,
not a causal estimate. At forecast horizon $h$, the head is
\begin{equation}
  \Delta\hat y_{\mathrm{ex},h}=-A_hR_hS,
  \label{eq:exercise_head}
\end{equation}
with
\begin{align}
  A_h &= \frac{\left[k * \operatorname{ReLU}
    (\Delta\mathrm{HR}-15)\right]_h}{A_{\mathrm{ref}}},
    \label{eq:exercise_activation}\\
  R_h &= \operatorname{ReLU}\left(\hat y_{\mathrm{base},h}-105\right),
    \label{eq:exercise_ramp}\\
  S &= \frac{16.7}{160-105}=0.303636.
    \label{eq:exercise_gain}
\end{align}
Here $A_h$ is an HR-intensity activation with a 15-bpm dead zone and a
causal short-lag kernel $k$; $R_h$ is a one-sided glucose-state safety ramp;
and $S$ is one cohort-wide scalar. There is no per-participant sensitivity
term. Participant exercise sensitivity was not recoverable from the state,
and the insulin-resistance hypothesis closed as a negative result: baseline
glucose was the only predictor of the subsequent drop.

The attempted gain optimization is itself informative. Training against the
scenario loss moved $S$ from $0.303636$ to $0.123238$. The scenario pinball
loss was $31{,}213$ times the magnitude-prior loss initially and $496$ times
the prior loss at the final epoch. This reproduces the known magnitude-muting
dynamic. It corroborates that the imposed magnitude cannot be learned through
this objective, so the drifted checkpoint was rejected and $S$ was frozen at
$0.303636$.

With frozen gain, the median effect in the $70$--$100$~mg/dL bin is
$0.0$~mg/dL because of safety gating. The response grows monotonically in
magnitude with $g_0$, with
$\operatorname{corr}(g_0,\Delta\hat y_{\mathrm{ex}})=-0.9402$ at the
45-minute peak. In the $160$--$200$~mg/dL bin, the median peak is
$-17.50$~mg/dL at 45 minutes and has begun to recover by 60 minutes
($-14.05$~mg/dL). Forecast MAE is unchanged at $10.0665$~mg/dL
($\Delta\mathrm{MAE}=0.0000$). The realized $17.50$~mg/dL peak can exceed
the $16.66$~mg/dL static anchor because the live $R_h$ trajectory can exceed
the static reference $R=55$; it was reported as observed and not retuned.

\begin{figure}[H]
  \centering
  \includegraphics[width=0.95\linewidth]{exercise/fig_exercise_path_b_s25.png}
  \caption{%
    \textbf{Path B imposed planning response.}
    Dose response across wearable-defined HR intensity and glucose-state
    gating at 60 minutes. Magnitude is imposed; AI-READI supplies timing,
    shape, and live glucose-state gating. This is a planning response, not a
    causal estimate. Comparison with the descriptive kernel is a wiring
    check, not independent validation, because that kernel supplied the
    magnitude anchor.%
  }
  \label{fig:exercise_pathb_doeffect}
\end{figure}

\begin{figure}[H]
  \centering
  \includegraphics[width=0.95\linewidth]{exercise/interp_cf_cards.png}
  \caption{%
    \textbf{Single-action Path B interpretation cards.}
    Rising-high, falling-low, and stable states show the frozen forecast,
    $q_{10}$--$q_{90}$ band, actual outcome, imposed do(exercise) overlay, and
    frozen-model driver attribution. No insulin action is shown because
    dosing is unavailable, and no carbohydrate action is shown because the
    meal counterfactual was not constructible. The falling-low state has
    $g_0=95$~mg/dL and imposed lowering of $0.0$~mg/dL, demonstrating safety
    gating. Per-window lowering uses the window's live $R_h$ and can exceed
    the bin-mean peak. The cards are descriptive and make no causal claim.%
  }
  \label{fig:exercise_pathb_cards}
\end{figure}

\subsection{Path C: An Identified, Deployable Exercise-Aware Hypoglycaemia Advisory}
\label{subsec:ex_pathc}

Path~C avoids a counterfactual. It is an identified predictive advisory built
on the frozen model and the measured signed-bias trajectory, evaluated against
the model's own forecasts and observed outcomes without a control arm. Its
scope is 3{,}813 moderate or vigorous non-insulin windows from 470 episodes,
including 105 observed hypoglycaemic outcomes.

\paragraph{Claim 1: ranking.}
The forecast-only risk-averse $q_{10}$ nadir ranks observed exercise-associated
hypoglycaemia substantially better than current glucose. Pooled AUROC is
$0.974$ without correction and $0.975$ with correction, versus $0.833$ for
current $g_0$. Corrected and uncorrected curves overlap because AUROC is
invariant to the horizon-specific rank-preserving shift. The ranking gain
comes from the streaming encoder reading impending exercise-like dynamics,
including rising HR. Median warning lead time is 30 minutes.

\paragraph{Claim 2: calibration at 70 mg/dL.}
Applying the signed-bias correction at the fixed 70~mg/dL threshold raises
sensitivity from $0.429$ to $0.762$ and rescues 35 near-miss hypoglycaemic
outcomes. The gain is concentrated in vigorous windows, where sensitivity
rises from $0.179$ to $0.846$ and the bias reaches $7.17$~mg/dL at 60 minutes;
26 of the 35 rescued outcomes are vigorous. The gain has a clear price:
false positives rise from 28 to 119, specificity falls from $0.992$ to
$0.968$, and positive predictive value falls from $0.616$ to $0.402$.
The correction can rescue only cases whose uncorrected nadir lies in the
near-miss band between 70~mg/dL and the strain-specific bias-adjusted upper
edge, 75.1~mg/dL for moderate and 77.2~mg/dL for vigorous windows. The 25
remaining missed outcomes are the frozen model's residual blind spot.

\begin{figure}[H]
  \centering
  \includegraphics[width=0.95\linewidth]{exercise/fig_s34_exercise_advisory.png}
  \caption{%
    \textbf{Exercise-aware hypoglycaemia advisory.}
    The ROC panel evaluates Claim~1, that the frozen forecast ranks observed
    hypoglycaemic outcomes better than current glucose. The rescued-zone
    panel evaluates Claim~2, that signed-bias correction improves calibration
    at the fixed 70~mg/dL threshold for near misses, with an accompanying
    false-alarm cost. This is a bias-corrected predictive warning on the
    frozen model, with no counterfactual and no causal claim.%
  }
  \label{fig:exercise_advisory}
\end{figure}

\subsection{Synthesis and Limitations}
\label{subsec:ex_synthesis}

The complete arc separates four contributions. First, a wearable-only
detector identifies ambulatory exercise-like episodes without glucose in the
detection path. Second, Stage~0 establishes that the counterfactual magnitude
is not identifiable on AI-READI because cross-patient matching offers overlap
without comparability, while within-patient matching offers comparability
without positivity. This is a transferable design finding, analogous to the
need for participant-disjoint evaluation splits. Third, Path~B provides an
honest imposed planning response whose timing, shape, and glucose-state gating
are data-derived but whose magnitude is not estimated. Fourth, Path~C provides
an identified advisory with a calibration benefit at a pre-specified threshold.

The detector remains wearable-defined rather than ground-truth-labelled and
is primarily sensitive to ambulatory activity. Meals, time of day, medication,
and selection can influence the descriptive kernel. Path~B inherits the
$n=118$ moderate-episode anchor and its imposed magnitude should not be read as
a causal estimate. Path~C has 105 hypoglycaemic outcomes, and its correction
improves threshold calibration rather than ranking. Insulin dosing is
unobserved, so insulin-treated participants are excluded from all causal scope;
the analyses here make no causal exercise claim even in the non-insulin cohort.

\begin{table}[H]
  \centering
  \caption{%
    \textbf{Wearable-defined exercise analysis summary.}
    The detector, residual audit, identifiability result, imposed planning
    response, and identified advisory are separate claims.%
  }
  \label{tab:exercise_arc_summary}
  \small
  \setlength{\tabcolsep}{3pt}
  \begin{tabular}{p{0.22\linewidth}p{0.31\linewidth}p{0.39\linewidth}}
    \toprule
    Item & Result & Interpretation \\
    \midrule
    Clean episodes & 2{,}171 across 788 participants & Primary wearable-defined set \\
    Strict recovery & 1{,}015 across 574 participants & Sensitivity subset, not a detection gate \\
    Refined onset & Median 30 min earlier & Backtracks to sustained step rise \\
    Negative controls & HR: 8.01 vs 0.27/$-0.46$ bpm; steps: 29.44 vs $-0.93$/$-0.23$ steps/min & Real onsets separate from shifted/random-awake controls \\
    Cadence & 1{,}859 slow; 223 walk; 81 sub-threshold; 6 brisk; 2 run-like & Slow-walk episodes dominate \\
    Glucose kernel & 16.66 mg/dL; nadir 28.31 min & Descriptive-associational anchor only \\
    Residual audit & No MAE gap; bias gap $+4.19$ vigorous and $+4.05$ moderate & Provisional residual headroom \\
    Stage 0 & Comparability and positivity cannot be achieved together & Magnitude not identifiable \\
    Path B & $S=0.303636$; $r=-0.9402$; $\Delta$MAE $=0$ & Coherent imposed planning response \\
    Path C & AUROC 0.975 vs 0.833 current $g_0$; 35 outcomes rescued & Identified advisory with a false-alarm cost \\
    \bottomrule
  \end{tabular}
\end{table}
